\documentclass[10pt]{article}
\usepackage[margin=20mm]{geometry}
\usepackage{amsmath,amssymb,mathtools,fontspec,unicode-math,xurl,hyperref}
\IfFontExistsTF{DejaVu Serif}{\setmainfont{DejaVu Serif}}{\setmainfont{Latin Modern Roman}}
\setmathfont{latinmodern-math.otf}
\setlength{\parindent}{0pt}\setlength{\parskip}{6pt}\sloppy
\title{Deletion–contraction, orientation and the specified Lefschetz class in braid variety cluster structures}
\author{AI-assisted technical review}
\date{}
\begin{document}
\maketitle
Abstract: The acquired publisher version of Galashin, Lam and Sherman-Bennett is examined through its Deodhar geometry, deletion–contraction argument, braid moves and cohomological application. A rank-one counterexample refutes Proposition 3.14, a used general criterion; a stronger factorial criterion adequate for the braid application is proved here. An explicit type A1 seed exhibits an orientation mismatch in the induction. Independently, a type A3 braid variety is a two-dimensional torus on which the specified form vanishes, refuting the stated Lefschetz conclusion even after correcting its dimension indices. These findings do not refute the coordinate-ring theorem. The assessment is Major on the affected derivations and application; adjudication and independent acceptance remain pending.\par
\section*{Identity, definitions and review scope}
The inspected work is P. Galashin, T. Lam and M. Sherman-Bennett, Braid variety cluster structures, II: general type, Inventiones mathematicae 243 (2025), 1079–1127, DOI 10.1007/s00222-025-01390-5 (https://doi.org/10.1007/s00222-025-01390-5). The acquired publisher version of record has 49 physical pages. Throughout, PDF pages are one-based; printed page numbers equal PDF page numbers plus 1078. The acquired arXiv:2301.07268v3 source is a separately identified artifact; equivalence with the publisher version is not assumed.\par
The working group in [1, PDF p. 2] is complex, simple and simply connected, with a fixed pinning and opposite Borel subgroups. A double word \(\beta=i_1\cdots i_m\) has letters in \(\pm I\) and Demazure product \(w_0\). Weighted flags are points of \(G/U_+\); strict relative position must be distinguished from relative position allowing a torus factor. The quotient \(\mathring R_\beta\) is asserted smooth, affine, factorial and irreducible, of complex dimension \[D=d(\beta)=m-\ell(w_0).\] The rightmost positive distinguished subexpression runs backwards from \(u_m=w_0\). Its solid positions form \(J_\beta\), with \(|J_\beta|=D\). The Deodhar torus has character variables \(x_e\), determined by their orders on the boundary divisors. A solid index is mutable when the corresponding almost-positive subexpression has initial value \(\mathrm{id}\). Frozen boundary divisors lie outside the open braid variety, and the frozen variables are units there.\par
Theorem 1.1 identifies \(\mathbb C[\mathring R_\beta]\) with the cluster algebra of the constructed seed. This statement is distinct from the general criterion in Proposition 3.14 and from the particular cohomological conclusion in Theorem 6.8. A seed on an affine variety, as defined on PDF p. 18, requires its initial torus to be an open dense subset. It does not already require every mutated variable to be regular on that variety.\par
The review covers the principal internal definitions, valuations, quotient and torus constructions, deletion–contraction hypotheses, induction ordering, selected braid calculations, folding dependencies and the exact Lefschetz assertion. Bounded native checks supplement the primary analysis. External algebraic-group, local-acyclicity and cohomological results remain identified dependencies where their proofs were not retrieved. Neither source indexing nor hash verification is counted as scientific validation.\par
\section*{Problem 1: the ring containment in Proposition 3.14 does not produce embedded cluster tori}
\subsection*{Claim and exact counterexample}
Proposition 3.14 [1, PDF pp. 21–22; printed pp. 1099–1100] states that if \(X\) is affine, factorial and irreducible, \(\Sigma\) is an integral sink-recurrent seed on \(X\), and \(\mathbb C[X]\subset U(\Sigma)\), then \(\mathbb C[X]=A(\Sigma)\). The proof first treats the cluster tori as embedded in \(X\), and then invokes a codimension-two complement. The stated containment does not supply those embeddings.\par
Equation reference: eq:rankone. Set \[X=\mathbb A^1\times\mathbb G_m,\qquad R=\mathbb C[x,y,y^{-1}],
 \qquad T=(\mathbb G_m)^2\subset X.\] Take \(x\) mutable, \(y\) frozen, \(d=(1,1)\), and exchange column \(\widetilde B=(0,1)^t\). A full skew matrix is \(\left(\begin{smallmatrix}0&-1\\1&0\end{smallmatrix}\right)\), with \(\omega=-\operatorname{dlog}x\wedge\operatorname{dlog}y\). Mutation gives \[
 xt=1+y,\qquad t=x'=(1+y)/x .\] The ring \(R\) is a localization of a polynomial unique factorization domain, so every geometric hypothesis on \(X\) holds. The seed is integral. Its mutable graph has one vertex and no arrows, hence it is sink-recurrent. Its exchange rows even span \(\mathbb Z\), so it is really full rank. The special frozen-row condition in the standing Assumption 3.11 also holds: the frozen row is \(+1\) and there is no other mutable column. Thus neither that standing assumption nor full rank eliminates the example.\par
There are only two alternating rank-one clusters, since mutation returns \(x\). Consequently \[U(\Sigma)=
 \mathbb C[x^{\pm1},y^{\pm1}]\cap\mathbb C[t^{\pm1},y^{\pm1}]
 \quad\hbox{inside }\mathbb C(x,y).\] Every generator of \(R\) belongs to both Laurent rings: in the second one, \(x=(1+y)/t\). Therefore \(R\subset U(\Sigma)\), exactly as required.\par
Nevertheless, \(t\) is a cluster generator and is not regular on \(X\). At the admissible point \((x,y)=(0,1)\), a regular extension of [eq:rankone] (\#eq:rankone) would imply \(0=2\). Equivalently, \(\operatorname{ord}_{\{x=0\}}t=-1\). Hence \(R\ne A(\Sigma)\). This is a counterexample to the proposition, with all its hypotheses verified, rather than merely an unfilled proof step.\par
The induced mutated-torus morphism is \[\operatorname{Spec}\mathbb C[t^{\pm1},y^{\pm1}]\longrightarrow X,\qquad
 (t,y)\longmapsto((1+y)/t,y).\] It sends every \((t,-1)\), \(t\ne0\), to \((0,-1)\). It is therefore not an open immersion. Together with the initial torus it still misses \(\{x=0,\ y\ne-1\}\), a codimension-one set. Exact reviewer-written arithmetic independently checked two distinct preimages and the nonzero numerator at \((0,1)\); the ring argument above establishes the general discrepancy. See the saved `aux\_induction/exact\_counterexample\_checks.json`.\par
\subsection*{Dependency and a sufficient repair}
Theorem 3.13 explicitly ends by invoking Proposition 3.14 [1, PDF p. 21]. This is a used dependency in the central deletion–contraction route to Theorem 1.1. In addition, Theorem 3.13 assumes only normality, whereas Proposition 3.14 assumes factoriality. The counterexample does not satisfy every stronger deletion–contraction hypothesis of Theorem 3.13, and is not claimed to refute that theorem itself.\par
The following replacement suffices for the actual factorial braid setting.\par
Proposition 1.  Let \(R=\mathbb C[X]\) be factorial and let \(\Sigma\) be integral, full rank and sink-recurrent on \(X\). Assume \(R\subset U(\Sigma)\), all initial variables and all one-step mutated variables belong to \(R\), and every frozen variable is a unit of \(R\). Then \(R=A(\Sigma)\).\par
Proof. Retain the paper’s local-acyclicity consequence \(A=U\) for sink-recurrent seeds. Every unit of \(A\) is a unit of the initial Laurent ring, hence a constant Laurent monomial. On mutating at a mutable index \(i\), the exchange polynomial \(P_i=M_++M_-\) is not a Laurent unit in the remaining variables: full rank excludes a zero exchange column, so the two monomials are distinct. Requiring both the unit and its inverse to remain Laurent forces its \(x_i\)-exponent to be zero. Thus units of \(A\) are precisely constant frozen Laurent monomials.\par
If an initial mutable \(x_i=ab\) in \(A\), the two factors are Laurent monomials on the initial torus. Mutation at \(j\ne i\) forces each factor’s \(x_j\)-exponent to vanish. Mutation at \(i\) forces its two exponents to be nonnegative and sum to one. One factor is therefore a frozen Laurent unit. The same proof applies to each adjacent seed, since full rank is mutation invariant. Initial mutable variables are pairwise nonassociate; \(x_i\) and \(x_i'\) are also nonassociate, since their exchange polynomial cannot equal a frozen monomial times \(x_i^2\).\par
These irreducibility and nonassociation assertions descend from \(A\) to \(R\): the functions are regular by hypothesis, and all units of \(A\) just identified are already units of \(R\). Factoriality makes the relevant irreducibles prime. For each initial or adjacent cluster, the inclusion \(R\subset U\) and regularity of its variables identify the localization of \(R\) at those variables with its Laurent ring. These are now actual principal-open tori in \(X\).\par
Every height-one prime is covered. If it contains an initial mutable variable \(x_i\), it cannot contain another initial mutable \(x_j\) or the nonassociate prime \(x_i'\); frozen variables are units. Thus the adjacent \(i\)-torus contains its generic point. The union of these tori has complement of codimension at least two. Normality gives \(U\subset R\), and hence \(R=U=A\). ◻\par
For braid varieties, factoriality is supplied by Proposition 2.5, and initial regularity and frozen units by Corollary 2.24. In the preceding argument of Theorem 3.13, adjacent regularity can be made explicit. The complement \(V=\{x_s=0\}\) is irreducible by its product description. For \(j\ne s\), the denominator \(x_j\) is nonzero at the generic point of \(V\), so \(x_j'=P_j/x_j\) has no pole there; it is already regular on the open deletion stratum. Normality extends it across \(X\). The special \(x_s'\) is regular by hypothesis. The product description also supplies the missing inverse on the exchange locus: the other torus coordinates specify the contraction point and \(x_s'\) specifies its affine-line coordinate.\par
This is a verified sufficient repair for the factorial application, conditional on the consistently oriented deletion–contraction hypotheses discussed next and on the cited \(A=U\) result. It does not independently settle the general normal, nonfactorial scope of Theorem 3.13. The finding is Major because a false criterion is explicitly used in a central derivation; a short stronger replacement does not erase that dependency.\par
\section*{Problem 2: the induction uses a source as a sink under the printed convention}
The graph convention on PDF p. 18 is \(i\to j\) when \(B_{ij}>0\). The sink criterion freezes a sink together with its incoming neighbors. Theorem 4.10 [1, PDF pp. 31–33; printed pp. 1109–1111] applies this construction to the leading mutable crossing of \(\beta=ii\beta'\).\par
Equation reference: eq:orientation. Take type \(A_1\), \(G=SL_2\), and \(\beta=1111\). Its Demazure product is the longest reflection. Writing \(\mathrm{id}=0\), \(s=1\), the positive distinguished sequence is \[(u_0,\ldots,u_4)=(0,0,0,0,1),\qquad J_\beta=\{1,2,3\}.\] The almost-positive sequences are \[(1,0,0,0,1),\quad(0,1,0,0,1),\quad(0,0,1,0,1).\] Hence vertex \(1\) is frozen and \(2,3\) are mutable. The vanishing-order criterion and simple diagonal zeros give, for the unique simple-root minors, \[\Delta_0=x_1,\quad\Delta_1=x_2,\quad\Delta_2=x_3,\quad\Delta_3=1.\] Using the internally consistent half-sum normalization of (2.26), discussed below, \(L_c=\operatorname{dlog}\Delta_c\). Equation (2.27) therefore gives \[
 \omega_\beta=\operatorname{dlog}x_1\wedge\operatorname{dlog}x_2+
                 \operatorname{dlog}x_2\wedge\operatorname{dlog}x_3,\qquad
 B=\begin{pmatrix}0&1&0\\-1&0&1\\0&-1&0\end{pmatrix}.\] Thus the mutable graph is \(2\to3\): vertex \(2\) is a source, not a sink. Its incoming-neighbor set is empty. Freezing \(\{2\}\cup N_2^{\mathrm{in}}\) leaves vertex \(3\) mutable, whereas \(\beta'=11\) has one frozen variable and no mutable variable. The contraction identification used in the proof therefore has the wrong mutable rank under the literal convention.\par
The exchange relation gives the geometric reason: \[x_2x_2'=x_1+x_3.\] On the contraction divisor \(x_2=0\), one has \(x_3=-x_1\), a unit. It is the outgoing neighbor \(3\) that must become frozen. The distinguished recurrences, skew matrix and neighbor sets were independently computed exactly in `aux\_induction/a1\_orientation\_check.json`.\par
A consistent correction reverses the declared graph convention globally, or uses the source/outgoing-neighbor version throughout the recurrence and application. The exchange matrix and exchange algebra need not change: reversing the quiver exchanges the two monomials in the symmetric exchange sum, and the corresponding local-acyclicity argument is mirrored. This removes the demonstrated local obstruction. It does not turn the literal contraction statement into a valid statement without that correction.\par
This is a substantive convention error at a used central induction step, with an explicit rank discrepancy, and is assessed Major for that derivation. It is not a counterexample to the resulting cluster-coordinate-ring identification. The correction is separate from the extra regularity hypotheses needed in Problem 1.\par
\section*{Problem 3: the specified class in Theorem 6.8 can vanish on a two-dimensional torus}
\subsection*{Exact source assertion and the index discrepancy}
Equation reference: eq:printedlefschetz. Section 6.4 [1, PDF p. 44; printed p. 1122] explicitly sets \(d=d(\beta)\), the full complex dimension. For even \(d\), it defines the curious Lefschetz property with respect to the particular class \([\omega_\beta]\) by \[
 [\omega_\beta]^{d-p}:
 H^{p+s,(p,p)}(\mathring R_\beta,\mathbb C)
 \longrightarrow H^{2d-p+s,(2d-p,2d-p)}(\mathring R_\beta,\mathbb C).\] Theorem 6.8 on the next page asserts this property for every even-dimensional double braid variety, and asserts the associated curious Poincare symmetry also in odd dimensions. The source images confirm both the dimension definition and the specified class; neither is an extraction artifact.\par
For a complex two-dimensional torus, [eq:printedlefschetz] (\#eq:printedlefschetz) with \(p=s=0\) would map the nonzero \(H^0\) isomorphically to \(H^4=0\). The dimension center has been doubled. With full dimension denoted \(D\), the expected even-dimensional indexing would instead use exponent \(D/2-p\) and target \[H^{D-p+s,(D-p,D-p)}.\] The printed Poincare symmetry has the same center discrepancy. Correcting those indices alone, however, does not repair the specified-class theorem.\par
\subsection*{An admissible braid with \(\omega_\beta=0\)}
Let \[G=SL_4,\qquad
 \beta=(1,3,1,2,1,3,2,1)=13\cdot121321 .\] The last six letters form a reduced expression of \(w_0\): right multiplication gives \[1234,\ 2134,\ 2314,\ 3214,\ 3241,\ 3421,\ 4321 .\] Thus \(\delta(\beta)=w_0\), \(m=8\), \(\ell(w_0)=6\) and \(D=2\), within every group and word assumption of the theorem.\par
The backwards distinguished recursion uses all six suffix letters and has \(u_2=u_1=u_0=\mathrm{id}\). Thus \(J_\beta=\{1,2\}\). At the first solid position, the almost-positive initial value is \(s_1\). At the second it is \(s_3\): indeed \(\min(s_3,s_3s_1)=s_3\), since \(1\) and \(3\) commute. Both solid positions are frozen. Proposition 2.19 consequently excludes every torus-boundary divisor from the open variety, giving \[\mathring R_\beta=T_\beta\cong(\mathbb C^\times)^2 .\]\par
Equation reference: eq:zeroform. Lemma 2.8 gives \[h_2^+=1,\qquad
 h_1^+=\alpha_3^\vee(t_2),\qquad
 h_0^+=\alpha_1^\vee(t_1)\alpha_3^\vee(t_2),\] and all later \(h_c^+=1\). The simple roots \(1\) and \(3\) are orthogonal: \(a_{13}=a_{31}=0\). Hence in the first local term of (2.27), \(L_{1,1}=0\); in the second, \(L_{2,3}=0\). Every hollow contribution vanishes by Corollary 2.27. Therefore \[
 \omega_\beta=0 .\] This conclusion is unaffected by the factor-of-two convention in (2.26).\par
Now \(H^0((\mathbb C^\times)^2,\mathbb C)\cong\mathbb C\) and \(H^{2,(2,2)}((\mathbb C^\times)^2,\mathbb C)\cong\mathbb C\). Even with the corrected dimension center, multiplication by \([\omega_\beta]\) from the former to the latter is zero and is not an isomorphism. The counterexample refutes the theorem with its specified class, beyond its indexing discrepancy.\par
Two bounded native checks independently reconstructed the Weyl permutations, frozen tests and zero wedge coefficients. The executed integer/rational calculations are saved in `aux\_geometry/low\_rank\_checks.json` and `aux\_folding/lefschetz\_counterexample.json`. They implement reviewer-written checks of the definitions, not author code.\par
\subsection*{Repair boundary and surviving conclusions}
On this particular torus, \(\eta=\operatorname{dlog}t_1\wedge\operatorname{dlog}t_2\) has a nonzero cohomology class and supplies the corrected-dimensional map \(H^0\to H^2\). Replacing a prescribed possibly degenerate form by a suitable nondegenerate one is a substantive change of statement. No general existence or compatibility theorem for such a replacement has been independently verified here. The cited external Lefschetz source has not been used as authority for the false assertion about an arbitrary specified \(\omega_\beta\).\par
The main coordinate-ring theorem survives this example: the all-frozen rank-zero seed gives exactly the Laurent ring \(\mathbb C[t_1^{\pm1},t_2^{\pm1}]\). Mixed Tate cohomology also holds for the torus. Corrected-center Poincare symmetry is not refuted by [eq:zeroform] (\#eq:zeroform); the printed-center version is refuted by the degree calculation. The Major assessment concerns the explicit cohomological application, not a claim that the entire cluster construction fails.\par
\section*{Problem 4: one folded mutation word does not project from its listed lift}
Table 1 [1, PDF p. 41; printed p. 1119], panels (b) and (f), gives the \(G_2\) case with hollow expanded positions \([8,11]\). The listed lifted word is \[\mu_{\mathrm{lift}}=\mu_{(5,6,7,12,5,6,7)},\] whereas the corresponding folded word is printed as \(\mu_{\mathrm{fold}}=\mu_{(6,3,6)}\). The expanded word and orbit map in the table caption give \[\lambda^{-1}(3)=\{5,6,7\},\qquad \lambda^{-1}(6)=\{12\}.\] Consequently this lift consists of orbit mutations \((3,6,3)\), not \((6,3,6)\). The assertion immediately before Proposition 6.4, and the literal use in Corollary 6.5 (PDF p. 42), require the two descriptions to correspond.\par
The discrepancy is not removed by reversing composition: both orbit words are palindromes. Commuting the three mutations inside the first orbit likewise does not change its order relative to the second orbit. Nor are the two folded mutation maps identical. With variables \((x_3,x_6)=(2,5)\), trivial coefficients and \[B=\begin{pmatrix}0&-3\\1&0\end{pmatrix},\] the two final ordered pairs are respectively \[\mu_{(3,6,3)}:(2,5)\mapsto(11/5,28/5),\qquad
 \mu_{(6,3,6)}:(2,5)\mapsto(7/5,52/25).\] These follow directly from \(x_3'=(1+x_6)/x_3\) and \(x_6'=(1+x_3^3)/x_6\), with matrix mutation between steps; all denominators are nonzero. The exact rational calculation and source transcription are preserved in `aux\_folding/table\_resolution.json`.\par
Changing this row of \(\mu_{\mathrm{fold}}\) to \((3,6,3)\) restores the projection from the listed lift. The independent principal-coefficient check of that lift against the braid sequence passed. This is assessed Minor as a demonstrated local table defect: the valid listed lift supplies counterevidence against concluding that the long-braid equivalence itself fails. The check does not certify a complete revised Table 1.\par
\section*{Positive checks, alternative readings and bounded coverage}
\subsection*{Geometry and low-rank boundaries}
The opposite weighted-flag pair has trivial stabilizer, but freeness alone would not prove affineness of a unipotent quotient. The proof has the stronger regular opposite-cell slice: after fixing \(X_0\), write the opposite flag \(Y_0=u w_0hU_+\), and use the unique \(u\in U_+\) to obtain the slice. The isolated \(Y_m\) in the slice sentence of Proposition 2.5 (PDF p. 7) must be \(Y_0\). For \(\beta=-i\) in rank one, \(Y_m=X_m=X_0\), making the literal index impossible, while the corrected slice works. The surrounding argument already identifies \(Y_0\); this indexing correction is recorded without a separate substantive Problem.\par
The chamber minor near a boundary has the form \((p+t)\Delta_{\omega_i,\omega_i}(h)\), with the second factor nonzero. This gives a simple diagonal zero and an integral unitriangular valuation matrix. The prescribed divisor orders consequently determine a character basis, rather than only a rational basis. Normality extends functions with no codimension-one poles. Frozen divisors are outside the open variety, explaining invertibility there. The general order criterion distinguishes zero from positive order; a universal \(0/1\) assertion was not imported from the separately stated \(SL_n\) case.\par
For \(SL_2\), direct determinant calculations give \[R_i\cong\{\mathrm{pt}\},\qquad R_{ii}\cong\mathbb C^\times,\qquad
 R_{iii}\cong\{(a,d)\in\mathbb C^2:1+ad\ne0\}.\] On the last surface, \(h=(1+ad)^{-1}\), \(x_1=-h,\ x_2=-ha,\ x_2'=-d\), and \(x_2x_2'=1+x_1\). Exact rational evaluations included the mutable boundary. The analytic determinant derivation, rather than those finitely many evaluations alone, establishes the formula.\par
Equation (2.26), PDF p. 16, equates a full logarithmic derivative with a half-sum. Later (4.6) and (4.15) explicitly use the half-root convention, which gives the rank-one exchange just checked. Taking the first expression literally instead multiplies the two-form by four and produces the wrong exchange exponent. The report uses the consistent half-sum throughout; the supplied alternative reading is therefore resolved by later source formulas, rather than promoted into a claim that the intended seed fails.\par
\subsection*{Moves, induction and external dependencies}
In the solid-special move, the minor identity reduces to \(A=B+1\), so \(\operatorname{dlog}A\wedge\operatorname{dlog}B=0\). In the nonspecial case, the difference of local forms is a multiple of \(d_j a'-d_i a=0\), by the invariant root pairing. The rank-two polynomial identity \[z_1(a)z_2(b)z_1(c)=z_2(c)z_1(ac-b)z_2(a)\] has both sides equal to \(\left(\begin{smallmatrix}ac-b&-a&1\\c&-1&0\\1&0&0\end{smallmatrix}\right)\); this directly supports the change of parameters in Proposition 4.9. These are analytic checks. They are not reported as executions of the author’s implementation.\par
Theorem 4.2 is conditional on an already established cluster-variety presentation. Its use of irreducibility in the final part of the move argument is therefore not itself circular. For a nonreduced positive word, the Coxeter exchange condition and braid equivalence expose a repeated letter; the deletion and contraction calls shorten the word. When the second leading crossing is hollow, the proof uses a product with \(\mathbb C^\times\) and only the first shorter word. It does not require an undefined empty-suffix variety in the minimal rank-one case. When both are solid, the genuine contraction applies.\par
The really-full-rank extension is an integral pivot-row calculation, using the special frozen coefficient \(\pm1\). Quasi-equivalence changes the character basis by an integral unimodular matrix, preserving the row lattice. Proposition 3.4’s unqualified irreducibility wording needs care for frozen units and zero exchange columns: in rank one with zero column, \(xx'=2\) makes \(x\) a unit over \(\mathbb C\). The actual full-rank range excludes this boundary case, and Proposition Proposition 1 (\#prop:repair) proves the required irreducibility there.\par
The proof ordering through Sections 5–7 was checked explicitly. The simply-laced main theorem precedes its use in folding; Proposition 7.2 assumes the main results, and Remark 7.3 explains that only its already established simply-laced instance is used in the multiply-laced proof. Its recursive argument increases the relevant Bruhat lengths until a terminal value, providing termination. This removes a superficial circularity objection. It does not independently re-prove every external generalized-minor, folding or local-acyclicity result.\par
\subsection*{Folding computations and a remaining labeling limitation}
The folding check distinguished orbit invariance from admissibility, retained the restriction to coefficients coming from the fixed-point embedding, and tracked frozen-variable rescaling separately from mutation. The reduction in Proposition 6.4 to framed restricted quivers was inspected in the source, including the use of the already proved simply-laced algorithm.\par
A reviewer-written exact calculation reconstructed the six nontrivial restricted-quiver cases displayed in Table 2: one \(B_2/C_2\) case and five \(G_2\) cases. It tracked principal-coefficient \(F\)-polynomials, gradings and exchange matrices. Five printed braid/lift comparisons passed. For the all-solid \(G_2\) case, the implemented right-to-left convention agreed after interchanging output labels \(11,12\); the printed exact labeling did not agree. Inverting both printed permutations did not resolve that discrepancy. Alternative left-to-right expansions reached the declared 30,000-term bound and were not completed.\par
The last result remains a bounded labeling concern requiring evidence, not an established additional failure of the general theorem. Mutation equivalence up to mutable relabeling is explicitly allowed in Proposition 6.6. No universal Table 1 certification is claimed. Inputs, exact outputs, convention trials and the preserved computation limit are saved in `aux\_folding/table\_mutation\_checks.json` and `aux\_folding/table\_resolution.json`. These scientific computations are distinct from the manifest and hash audits.\par
\section*{Overall impact and state}
Problems 1 and 2 concern used dependencies of the coordinate-ring proof. The first is repaired for the actual factorial setting by added regularity and unit hypotheses; the second requires a consistent graph orientation. Their evidence does not show that Theorem 1.1 is false. Problem 4 is a local folded-word transcription mismatch with a verified projection correction. Problem 3 is a counterexample to the particular Lefschetz-class claim itself. Correcting its dimension center is necessary but insufficient; a general replacement class or extra hypothesis remains to be justified.\par
No author correction was independently acquired or adopted in this recovery. The report does not claim priority of discovery. Its demonstrated counterexamples and verified local repairs are distinguished from the unresolved normal-only criterion and any general replacement cohomological theorem.\par
Affected claim | Assessment | Repair verification | Adoption\par
Proposition 3.14 and its use | Major | Stronger factorial criterion proved; normal-only generality unresolved | Unknown\par
Induction orientation | Major | Explicit source/outgoing correction and local rank check | Unknown\par
Specified Lefschetz class | Major | Index correction insufficient; alternative form checked only on the counterexample | Unknown\par
Folded mutation word | Minor | Corrected orbit projection and listed lift checked | Unknown\par
Scientific adjudication is pending; `human\_reviews=[]`. This is an initial recovery review, not independent acceptance.\par
1 P. Galashin, T. Lam and M. Sherman-Bennett, Braid variety cluster structures, II: general type, Invent. Math. 243 (2025), 1079–1127, doi:10.1007/s00222-025-01390-5. Acquired publisher PDF, 49 physical pages. All numbered source claims in this report refer to this exact version.\par
\section*{Supporting checks}
The canonical publisher PDF SHA256 is\par
78b791f94f0d7e6abaca00d2acb86eb37de7
9eb62021cd9f8cb445fa7dcc9610\par
(The two lines concatenate without spaces.) The exact manifest, source, report, compiler, dependency, log and image hashes are stored in the per-paper verification records; hashes are derived from real bytes.\par
and the recovery token is\par
The previous preserved inputs contained no recoverable substantive scientific report for this paper. This review therefore supplies the necessary initial science and report, while retaining the prior source manifest and acquisition attribution.\par
The primary session performed source reading, counterexample analysis, repair integration and report preparation. These are AI support checks, not separate independent acceptance. The external controller reports the UI selection Personal/Chat/GPT-6/Pro; backend identity was not independently verified. No human review is inferred.\par
Actual source and final-report page views are bound separately from renders. The existing Tectonic compiler runs with `–untrusted` and `–only-cached`. Its execution records retain logs, intermediates, makefile dependencies and before/after support-file hashes; those inventories are conservative supersets. No author code, installation or shell escape is used.\par
\end{document}
